% =====================================================================
%  c23 —— 原文印刷页 182--191（PDF 第 183--192 页）
%  含：附录 D（D. Level k Theta Functions；D.1 定义、D.2 模变换律）
%      附录 E（E. Theta Functions And Holomorphic Line Bundles On Tori；
%      E.1 Heisenberg 群作为主 U(1) 丛、E.2 自然联络、
%      E.3 辛环面上的复结构、E.4 拉格朗日分解、
%      E.5 环面上复结构的空间：Siegel 上半平面）
%  编号公式：D.1--D.21、E.1--E.45 共 66 条；本片段无插图。
%  全部公式均逐条对照 150 dpi 页面图像转写（文本层仅供读散文）。
% =====================================================================

\section{第 $k$ 级 theta 函数} % 原文: D. Level k Theta Functions

\begin{keytakeaways}
  \item 定义：第 $k$ 级 theta 函数 $\Theta_{\mu,\kappa}(z,\tau)$ 是 $q^{\ell^2/(4\kappa)}y^{\ell}$
    在 $\ell\equiv\mu\ (\mathrm{mod}\ 2\kappa)$ 上的求和；$\mu$ 只需看成 $\mathbb{Z}/2\kappa\mathbb{Z}$
    的元素，因而同一模类只给出一个函数。
  \item 拟周期性：沿 $z\mapsto z+\nu$ 严格不变，沿 $z\mapsto z+\nu\tau$ 则多出乘性因子
    $e^{-2\pi i\kappa\nu^2\tau-4\pi i\kappa\nu z}$；该因子与 $\mu$ 无关，故所有模类「以同一方式」变换。
  \item 模变换：把 $\Theta_{\mu,\kappa}$ 除以 $\eta(\tau)$ 得到 $f_{\mu,k}$，其张成空间 $V_k$ 是
    $SL(2,\mathbb{Z})$（而非 $PSL(2,\mathbb{Z})$）的表示；按 $-1$ 的本征值分解为
    $\dim V_k^+=k+1$、$\dim V_k^-=k-1$ 的两块。
  \item 生成元的矩阵元：$\rho(S)_{\mu\nu}=\frac{1}{\sqrt{2k}}e^{2\pi i\mu\nu/(2k)}$ 是 Gauss 和型相位，
    $\rho(T)$ 的相位 $\mu^2/(4k)-1/24$ 中的 $-1/24$ 正是中心荷为 $1$ 的 $\eta$ 因子的残留。
  \item 一般元素的变换律 (D.16) 需要两个数论输入：Rademacher 函数 $\Phi(A)$（由 Dedekind 和给出，
    取整值）与连分数项数 $t$（只在模 $2$ 意义下良定义，故 $(-1)^t$ 良定义）。
\end{keytakeaways}

\subsection{定义} % 原文: D.1 Definition

这里汇集关于 theta 函数的若干事实。回忆
\begin{equation}
\begin{aligned}
  \Theta_{\mu,\kappa}(z,\tau) &:= \sum_{\ell=\mu\,\mathrm{mod}\,2\kappa} q^{\ell^2/(4\kappa)} y^{\ell}\\
  &= \sum_{n\in\mathbb{Z}} q^{\kappa(n+\mu/(2\kappa))^2} y^{(\mu+2\kappa n)}\\
  &= \sum_{n\in\mathbb{Z}} e^{i\pi\tau(2\kappa)(n+\mu/(2\kappa))^2+2\pi i(2\kappa)z(n+\mu/(2\kappa))}
\end{aligned}
\tag{D.1}
\label{eq:D.1}
\end{equation}
其中 $q=e^{2\pi i\tau}$，$y=e^{2\pi iz}$。这里 $\mu$ 是整数，$\kappa$ 是正半整数（即属于
$\frac12\mathbb{Z}_+$）。注意若把 $\mu$ 平移为 $\mu\to\mu+2\kappa s$，$s$ 为任意整数，则
$\Theta_{\mu,\kappa}$ 不变。人们常把 $\mu$ 取在基本域 $-\kappa<\mu\le\kappa$ 内，但一般说来应当把
$\mu$ 视为 $\mathbb{Z}/2\kappa\mathbb{Z}$ 的元素。作为 $z$ 的函数，这些函数是双拟周期的：
\begin{equation}
\begin{aligned}
  \Theta_{\mu,\kappa}(z+\nu,\tau) &= \Theta_{\mu,\kappa}(z,\tau)\\
  \Theta_{\mu,\kappa}(z+\nu\tau,\tau) &= e^{-2\pi i\kappa\nu^2\tau-4\pi i\kappa\nu z}\Theta_{\mu,\kappa}(z,\tau)
\end{aligned}
\tag{D.2}
\label{eq:D.2}
\end{equation}
这里 $\nu$ 是任意整数。注意对一切 $\mu\in\mathbb{Z}/2\kappa\mathbb{Z}$，theta 函数的变换方式都相同。
关于这些 theta 函数的几何意义，我们将在下文的
{\small\itshape ［原文标记：section ****（章节号待补，原稿占位符）］} 中作更多解释。
% [原文待办标记，原文 p.182：正文写的是「in section **** below」，星号是作者留给自己的章节号占位符]

% [Tutor Note: (D.1) 三行是同一个求和的三种写法，核对时最易错的是指数。第一行的 $\ell$ 跑遍与
% $\mu$ 同余模 $2\kappa$ 的整数，故 $\ell^2/(4\kappa)$ 与第二行的 $\kappa(n+\mu/(2\kappa))^2$ 相等
% 只须代入 $\ell=\mu+2\kappa n$：$(2\kappa)^2/(4\kappa)=\kappa$，这正是「$4\kappa=(2\kappa)^2/\kappa$」
% 那一步。第三行把 $q=e^{2\pi i\tau}$、$y=e^{2\pi iz}$ 代回，于是 $\tau$ 前多出因子 $2\kappa$、
% $z$ 前也多出因子 $2\kappa$——这两处因子最容易漏掉其中一个。]

\subsection{第 $k$ 级 theta 函数的模变换律} % 原文: D.2 Modular Transformation Laws For Level k Theta Functions
{\small\itshape ［原文标记：此处需要把约定 $k$ 改成 $\kappa$］}
% [原文待办标记，原文 p.182 右侧页边注：\ensuremath{\clubsuit} Need to change conventions k → κ here \ensuremath{\clubsuit}]

实际上，用下面这些函数来工作会更方便：
\begin{equation}
  f_{\mu,k}(z,\tau) := \frac{\Theta_{\mu,k}(z,\tau)}{\eta(\tau)}
  \tag{D.3}
  \label{eq:D.3}
\end{equation}
把这些函数的张成（span）记作 $V_k$。它将是 $SL(2,\mathbb{Z})$ 的一个表示（而\emph{不是}
$PSL(2,\mathbb{Z})$ 的表示），我们的目标就是确定这个表示。对应于 $-1\in SL(2,\mathbb{Z})$ 的作用的
对合 $z\to -z$ 可以对角化，从而把 $V_k$ 分解为偶子空间与奇子空间的直和
\begin{equation}
  V_k\cong V_k^+\oplus V_k^-
  \tag{D.4}
  \label{eq:D.4}
\end{equation}
于是 $\dim V_k^+=k+1$，$\dim V_k^-=k-1$。奇空间 $V_k^-$ 中一般第 $k$ 级 theta 函数的变换规则，
是 Lisa Jeffrey [31] 在一项令人印象深刻的计算中导出的。借助她论文中的想法，可以把结果推广为
$V_k$ 中全部 theta 函数的变换律。这个计算略有些非平凡，所以此处直接给出结果（因为它在别处似乎
还查不到）。

我们在 $(z,\tau)$ 上有一个 $SL(2,\mathbb{Z})$ 的\emph{左}作用：若
\begin{equation}
  A = \begin{pmatrix} a & b\\ c & d\end{pmatrix}
  \tag{D.5}
  \label{eq:D.5}
\end{equation}
则
\begin{equation}
  A\cdot(z,\tau) := \left(\frac{z}{c\tau+d},\ \frac{a\tau+b}{c\tau+d}\right)
  \tag{D.6}
  \label{eq:D.6}
\end{equation}
定义
\begin{equation}
  \phi_k(A;(z,\tau)) = \exp 2\pi i k\,\frac{cz^2}{c\tau+d}
  \tag{D.7}
  \label{eq:D.7}
\end{equation}
那么
\begin{equation}
  \phi_k(A_1;A_2\cdot(z,\tau))\,\phi_k(A_2;(z,\tau)) = \phi_k(A_1A_2;(z,\tau))
  \tag{D.8}
  \label{eq:D.8}
\end{equation}
现在在 $V_k$ 上定义一个\emph{左}作用 $SL(2,\mathbb{Z})$：
\begin{equation}
  (\rho(A)\cdot f)(z,\tau) := \left(\phi_k(A^{-1};(z,\tau))\right)^{-1} f(A^{-1}\cdot(z,\tau))
  \tag{D.9}
  \label{eq:D.9}
\end{equation}
于是，相对于基 $f_{\mu,k}$，我们有
\begin{equation}
  \rho(A)\cdot f_{\mu,m} := \sum_{\nu}\rho(A)_{\nu\mu}f_{\nu}
  \tag{D.10}
  \label{eq:D.10}
\end{equation}
其中再次把 $\mu,\nu$ 视为模 $2k$ 的整数。

若 $a=0$，则 $A=S$ 或 $S^{-1}$，其中我们定义
\begin{equation}
  S := \begin{pmatrix} 0 & -1\\ 1 & 0\end{pmatrix}
  \tag{D.11}
  \label{eq:D.11}
\end{equation}
若 $c=0$ 且 $a=1$，则 $A=T^m$，此时只须算出
\begin{equation}
  T := \begin{pmatrix} 1 & 1\\ 0 & 1\end{pmatrix}
  \tag{D.12}
  \label{eq:D.12}
\end{equation}
的矩阵就够了。

若 $c=0$ 且 $a=-1$，则只须算出 $A=-1$ 的作用。直接计算给出：
\begin{equation}
  \rho(-1)_{\mu\nu} = \delta_{\mu+\nu,0}
  \tag{D.13}
  \label{eq:D.13}
\end{equation}
这就是电荷共轭矩阵。
\begin{equation}
  \rho(S)_{\mu\nu} = \frac{1}{\sqrt{2k}}\exp\left(2\pi i\,\frac{\mu\nu}{2k}\right)
  \tag{D.14}
  \label{eq:D.14}
\end{equation}
\begin{equation}
  \rho(T)_{\mu\nu} = \delta_{\mu,\nu}\exp\left(2\pi i\left[\frac{\mu^2}{4k}-\frac{1}{24}\right]\right)
  \tag{D.15}
  \label{eq:D.15}
\end{equation}
若 $ac\ne 0$，则
\begin{equation}
  \rho(A)_{\mu\nu} = \frac{1}{\sqrt{2k|c|}}
  e^{\frac{i\pi}{2k}\frac{b}{a}\nu^2}
  e^{-\frac{2\pi i}{24}\Phi(A)}
  \sum_{\gamma\,\mathrm{mod}(2kc),\ \gamma=\mu\,\mathrm{mod}\,2k}
  \exp\left[i\pi\frac{a}{2kc}\left(\gamma+(-1)^{t+1}\frac{\nu}{a}\right)^2\right]
  \tag{D.16}
  \label{eq:D.16}
\end{equation}
这里 $\Phi(A)$ 是 Rademacher 函数，$t$ 是 $A$ 的任一连分数展开中的项数。虽然 $A$ 可以容许许多种
连分数展开，我们断言项数模 $2$ 是良定义的，因此 $(-1)^t$ 是 $A$ 的良定义函数。

% [Tutor Note: (D.10) 与 (D.13)--(D.16) 的指标次序不同，极易混淆。(D.10) 里写的是
% $\rho(A)_{\nu\mu}$，即第一个下标是行、第二个下标是列，而 $\sum_\nu$ 收缩的是列指标，这正是
% 「左作用」的样子；(D.13)--(D.16) 则直接给出矩阵元 $\rho(A)_{\mu\nu}$，两者相差一次转置。
% 若照抄时把 (D.10) 也写成 $\rho(A)_{\mu\nu}$，后面 (D.17) 的检验（$\gamma\to-\gamma$ 换元）就会
% 给出错误的对称性；另外 (D.10) 里基函数的下标原作 $f_{\mu,m}$，与全片的 $f_{\mu,k}$ 不一致，
% 疑为笔误，此处照原样转写。]

% [Tutor Note: (D.16) 的两处相位要分开看：$e^{-\frac{2\pi i}{24}\Phi(A)}$ 与 (D.15) 里方括号中的
% $-1/24$ 同源，都来自 $\Theta/\eta$ 中 $\eta(\tau)$ 的中心荷 $c=1$ 相位（$e^{2\pi i/24}$ 的十二次
% 幂为 $1$，故它只在模 $24$ 意义下确定）；而 $e^{\frac{i\pi}{2k}\frac{b}{a}\nu^2}$ 则是
% 「$b/a$ 型」二次相位的高斯和。核对 (D.16) 的关键是：把 $A=S$（$a=0$）代入会发散，所以 (D.16)
% 只在 $ac\ne0$ 时成立，$a=0$ 与 $c=0$ 分别由 (D.14)、(D.15) 单独给出。$(-1)^{t+1}$ 中的 $t$
% 必须取自同一次连分数展开，因为 $t$ 本身只在模 $2$ 下良定义——这一点正是下一段所强调的。]

（见 {\small\itshape ［原文标记：section ??（交叉引用的章节号未定）］}。）正是这个符号构成了与
Lisa Jeffrey 就奇 theta 函数所得表达式之间最重要的差别。
% [原文待办标记，原文 p.184：正文写的是「(See section ??.)」，问号是未定的 \ref 占位符]

注意
\begin{equation}
  \rho(A)_{\mu,\nu} = \rho(A)_{-\mu,-\nu}
  \tag{D.17}
  \label{eq:D.17}
\end{equation}
对一切 $A\in SL(2,\mathbb{Z})$ 成立。这必定如此，因为代表 $-1$ 的电荷共轭必须与所有矩阵对易。
就 (D.16) 的情形而言，这也可以通过把求和变量改为 $\gamma\to-\gamma$ 来验证。

剩下要定义的就是 $t$ 与 Rademacher 函数。

Rademacher 函数的定义与 [31] 中完全一样。函数 $\Phi:SL(2,\mathbb{Z})\to\mathbb{Z}$ 由下式定义
\begin{equation}
  \Phi(A) :=
  \begin{cases}
    \dfrac{a+d}{c}-12\operatorname{sign}(c)\,s(d,|c|) & c\ne 0\\[2mm]
    \dfrac{b}{d} & c=0
  \end{cases}
  \tag{D.18}
  \label{eq:D.18}
\end{equation}
其中 $s(d,c)$ 是 Dedekind 和，对 $c>0$ 定义为
\begin{equation}
  s(d,c) = \frac{1}{4c}\sum_{k=1}^{c-1}\cot\frac{\pi k}{c}\cot\frac{\pi dk}{c}
  \tag{D.19}
  \label{eq:D.19}
\end{equation}
（当 $d\ne 0$），并约定 $s(0,1)=0$。$\Phi$ 取整值这一点远非显然，但它可由下面这个值得注意的公式
推出
\begin{equation}
  \Phi(A_{12}) = \Phi(A_1)+\Phi(A_2)-3\operatorname{sign}(c_1c_2c_{12})
  \tag{D.20}
  \label{eq:D.20}
\end{equation}
其中 $A_{12}=A_1A_2$。

利用连分数理论可以证明，任何 $A\in SL(2,\mathbb{Z})$ 都具有如下形式
\begin{equation}
  A = T^{m_t}ST^{m_{t-1}}S\cdots T^{m_1}S
  \tag{D.21}
  \label{eq:D.21}
\end{equation}
其中——这一点很重要——所有 $m_i\ne 0$。由于群中的关系式，整数 $t$ 只在模 $2$ 意义下良定义。

% [Tutor Note: 作者写「$\Phi$ 取整值……远非显然」，跳过的是 Dedekind 和的三项关系
% （Dedekind reciprocity）：正是这条经典互反公式保证 (D.20) 中的组合在 $A_{12}=A_1A_2$ 处取整值，
% 进而由一个归纳论证给出 $\Phi(A)\in\mathbb{Z}$。要补齐这一步，可查 Rademacher 关于 Dedekind 和的
% 原始工作，或 Zagier 关于 Dedekind 和互反公式的综述；[31] 中只按定义引用而未复述证明。]

\section{环面上的 theta 函数与全纯线丛} % 原文: E. Theta Functions And Holomorphic Line Bundles On Tori

\begin{keytakeaways}
  \item 统一的视角：环面上的 theta 函数是某个主 $U(1)$ 丛（Heisenberg 群的中心扩张）的全纯截面，
    而模变换律来自丛的构造本身，与截面的具体选法无关。
  \item Heisenberg 群：整值辛形式 $\Omega:\Lambda\times\Lambda\to\mathbb{Z}$ 诱导中心扩张
    $0\to U(1)\to\tilde V\to V\to 0$；在 $\Lambda$ 上取分裂 $\epsilon$ 后商掉 $\Lambda_\epsilon$，
    即得环面上的主 $U(1)$ 丛 $(V\times U(1))/\Lambda$，其量子化维度为 $\dim\mathcal{H}=e_1\cdots e_n$，
    相应的线丛就是预量子线丛。
  \item 自然联络：平行移动由 $V$ 中的平移给出，曲率恰为 $\Omega$；第一 Chern 类由
    $(e_1,\dots,e_n)\in H^2(V/\Lambda;\mathbb{Z})$ 给出，闭路的和乐为
    $e^{2\pi i\Omega(\lambda,w_0)}\epsilon_\lambda^{-1}$。
  \item 复结构：取与 $\Omega$ 相容的 $J$，则 $g(v,w)=\Omega(Jv,w)$ 对称；周期矩阵
    $\tau_{IJ}$ 落在 Siegel 上半平面 $\mathcal{H}_N$ 内，且 $\mathcal{H}_N\cong Sp(2N,\mathbb{R})/U(N)$。
  \item 模空间：$\mathbb{R}^{2N}/\mathbb{Z}^{2N}$ 上复结构的空间是
    $Sp(2N,\mathbb{Z})\backslash\mathcal{H}_N$；$Sp(2N,\mathbb{Z})$ 由三类生成元给出，这正是 theta
    函数具有如此模变换性质的根源。
\end{keytakeaways}

在处理阿贝尔 Chern-Simons 理论时，对 theta 函数及其与环面上线丛截面的关系有一个一般的看法是很有用的。
这里把这套理论概述一遍。

若干数学处理可参见

\begin{enumerate}
  \item D. Mumford，\textit{Tata Lectures on Theta}，第 I 卷，p.202。那里对 theta 函数的变换
    有很好的描述。
  \item D. Mumford、M. Nori 与 P. Norman，\textit{Tata Lectures on Theta}，第 III 卷。他们把
    $\theta$ 解释为 Fock 空间表示中的一个特殊向量，也把 $\theta$ 解释为 Heisenberg 表示中的一个
    特殊矩阵元。
  \item 在代数几何框架内的另一种很好的描述可见 Griffiths 与 Harris，p.300 及其后。
  \item 专门讨论这一主题的一本很好的书是 Birkenhake 与 Lange，
  \item 这一主题在数以百计的物理学论文中从各种观点被处理过。我们用到了 [9] 的 pp.68--69 以及
    [20] 的第 6 节。
\end{enumerate}

\subsection{作为主 $U(1)$ 丛的 Heisenberg 群} % 原文: E.1 Heisenberg Groups As Principal U(1) Bundles

假设 $\Lambda$ 是一个秩为 $2n$ 的格，其上有整值辛形式：
\begin{equation}
  \Omega:\Lambda\times\Lambda\to\mathbb{Z}
  \tag{E.1}
  \label{eq:E.1}
\end{equation}
它诱导出实 $2n$ 维向量空间 $V=\Lambda\otimes\mathbb{R}$ 上的辛形式。利用 $\Omega$ 我们定义一个中心扩张
\begin{equation}
  0\to U(1)\to\tilde V\to V\to 0.
  \tag{E.2}
  \label{eq:E.2}
\end{equation}
可以把 $\tilde V$ 看成数对 $(v,z)\in V\times U(1)$ 的集合，其乘法由一个上闭链（cocycle）给出
\begin{equation}
  (v,z)(v',z') = (v+v',zz'c(v,v'))
  \tag{E.3}
  \label{eq:E.3}
\end{equation}
使得
\begin{equation}
  \frac{c(v,v')}{c(v',v)} = e^{2\pi i\Omega(v,v')}.
  \tag{E.4}
  \label{eq:E.4}
\end{equation}
{\small\itshape ［原文标记：若把约定改一改会更好，即把中心部分放在第一个因子上，于是元素写成
$(z,v)$ 而不是 $(v,z)$。这在接下来的全篇中都需要改。］}
% [原文待办标记，原文 p.185 右侧页边注：\ensuremath{\clubsuit} Would be better to change convention and put the central
% part as the first factor, so write elements as (z,v) rather than (v,z). That needs to be changed
% throughout the following. \ensuremath{\clubsuit}]

自然地取 $c(v,v')=e^{i\pi\Omega(v,v')}$，它对应于主 $U(1)$ 丛 $\tilde V\to V$ 的一个特殊平凡化。

现在可以构造环面 $V/\Lambda$ 上的主 $U(1)$ 丛了。为此我们把序列 (E.2) 在 $\Lambda$ 上分裂。也就是说，
选取一个函数 $\epsilon:\Lambda\to U(1)$，满足
\begin{equation}
  \epsilon_v\epsilon_{v'} = e^{-i\pi\Omega(v,v')}\epsilon_{v+v'}\qquad \forall v,v'\in\Lambda
  \tag{E.5}
  \label{eq:E.5}
\end{equation}
{\small\itshape ［原文标记：需要解释这样的分裂为何存在。］}
% [原文待办标记，原文 p.185 右侧页边注：\ensuremath{\clubsuit} Explain why such a splitting exists. \ensuremath{\clubsuit}]

现在 (E.5) 正是我们所需要的条件，使得 $v\to(v,\epsilon_v)$ 定义出同态 $\Lambda\to\tilde V$，即把
$\Lambda$ 嵌入群 $\tilde V$ 中。

把这个同态的像记作 $\Lambda_\epsilon$。现在取商 $\tilde V/\Lambda_\epsilon$ 就有意义了。这个空间是
$V/\Lambda$ 上的一个主 $U(1)$ 丛。

作为空间，我们的 $U(1)$ 丛可以写成 $(V\times U(1))/\Lambda$，其中的等价关系是
\begin{equation}
  (v,z)\sim(v,z)\cdot(\lambda,\epsilon_\lambda) = (v+\lambda,\epsilon_\lambda e^{i\pi\Omega(v,\lambda)}\cdot z)
  \tag{E.6}
  \label{eq:E.6}
\end{equation}

% [Tutor Note: 作者在 (E.5) 后让我们「相信分裂存在」，其实这一步可以显式写出：因 $\Omega$ 在
% $\Lambda$ 上取整值，$\Omega(\lambda,\lambda)\in\mathbb{Z}$，故 $\epsilon_\lambda:=e^{i\pi\Omega(\lambda,\lambda)/2}$
% 取值只可能是 $\pm1$，把它代入 (E.5) 并利用 $\Omega(v+v',v+v')=\Omega(v,v)+\Omega(v',v')+2\Omega(v,v')$
% 即得 $\epsilon_v\epsilon_{v'}\epsilon_{v+v'}^{-1}=e^{-i\pi\Omega(v,v')}$，条件恰好成立。也就是说，
% 「分裂存在」等价于 $\Omega$ 的整性；若 $\Omega$ 只是实值的，还可以差一个特征标，这正对应
% 下面第 2 条注记中的「二次细化」。]

\paragraph{注记}
\begin{enumerate}
  \item 在上述扩张中 $z\in U(1)$，所以我们构造的是主 $U(1)$ 丛。用这样的描述很容易得到相应的
    相伴丛。下文要描述环面上的全纯丛，届时这一点会很重要。只要在上述各式中让 $z$ 取值于
    $\mathbb{C}^*$ 或 $\mathbb{C}$，我们就能直接定义相伴的 $\mathbb{C}^*$ 或 $\mathbb{C}$ 丛！
    用 $\mathbb{C}^*$ 作的扩张会收缩到用 $U(1)$ 作的扩张上。
    {\small\itshape ［原文标记：这第一条注记位置不当，或者根本不必要。］}
    % [原文待办标记，原文 p.187 右侧页边注：\ensuremath{\clubsuit} This first remark is out of place or just not necessary. \ensuremath{\clubsuit}]
  \item 由于 $\Omega$ 在 $\Lambda\times\Lambda$ 上整值，双线性形式 $e^{i\pi\Omega(v,v')}$ 实际上是对称
    双线性形式。一般地，给定阿贝尔群 $A$ 上的对称形式
    \begin{equation}
      b:A\times A\to\mathbb{R}/\mathbb{Z}
      \tag{E.7}
      \label{eq:E.7}
    \end{equation}
    $b$ 的一个\emph{二次细化}（quadratic refinement）是函数 $q:A\to\mathbb{R}/\mathbb{Z}$，满足
    \begin{equation}
      q(x+y)-q(x)-q(y)+q(0)=b(x,y)
      \tag{E.8}
      \label{eq:E.8}
    \end{equation}
    注意我们可以把 $\epsilon$ 看作这个对称……
    （原文此处断句，句子在此中断；所缺的宾语即「对称形式 $b$ 的二次细化」，
    把 (E.5) 取对数并与 (E.8) 对照即可看出）
  \item 这给出了\emph{紧相空间}的一个漂亮例子。$T(V/\Lambda)$ 是辛流形。这类相空间的量子化应当
    给出一个有限维 Hilbert 空间，而在 WKB 近似下其维数应当是
    \begin{equation}
      \dim\mathcal{H} = \int_T\frac{\Omega^n}{n!}
      \tag{E.9}
      \label{eq:E.9}
    \end{equation}
    现在，$\Lambda$ 上任一整值反对称形式都可以看成一个矩阵，而在格的基变换
    $\Omega\to A^{tr}\Omega A$（$A\in GL(2n,\mathbb{Z})$）之下，我们可以把 $\Omega$ 化成如下形式：
    \begin{equation}
      \Omega = -ie_1\sigma^2\oplus -ie_2\sigma^2\oplus\cdots
      \tag{E.10}
      \label{eq:E.10}
    \end{equation}
    {\small\itshape ［原文标记：请给出证明。］}
    % [原文待办标记，原文 p.187 右侧页边注：\ensuremath{\clubsuit} Indicate the proof \ensuremath{\clubsuit}]
    其中 $e_I$，$I=1,\dots,n$ 是整数。我们可以要求 $e_1|e_2|\cdots|e_n$ 来把它们唯一确定。这就定义
    了所谓的\emph{辛基}。这一点后文还会再谈。此时 $\dim\mathcal{H}=e_1\cdots e_n$。在几何量子化中，
    与我们 $U(1)$ 丛相伴的线丛就是预量子线丛（pre-quantum line bundle）。
\end{enumerate}

\paragraph{练习}
说明主 $U(1)$ 丛还可以有如下另一种表示。由上闭链 $c((x,y),(x',y'))=xy'$ 构造 Heisenberg 表示
\begin{equation}
  0\to\mathbb{R}\to Heis(\mathbb{R}\times\mathbb{R})\to\mathbb{R}\times\mathbb{R}\to 0
  \tag{E.11}
  \label{eq:E.11}
\end{equation}
它有一个漂亮的 $3\times 3$ 矩阵表示
\begin{equation}
  \begin{pmatrix} 1 & x & z\\ 0 & 1 & y\\ 0 & 0 & 1\end{pmatrix}
  \tag{E.12}
  \label{eq:E.12}
\end{equation}
从拓扑上说这不过是 $\mathbb{R}^3$，但群乘法是非交换的。注意
\begin{equation}
  0\to\mathbb{Z}\to Heis(\mathbb{Z}\times\mathbb{Z})\to\mathbb{Z}\times\mathbb{Z}\to 0
  \tag{E.13}
  \label{eq:E.13}
\end{equation}
构成一个离散子群，只要取 $x,y,z$ 为整数即可。于是把 $\mathbb{R}^3/\mathbb{Z}^3$ 实现为
$Heis(\mathbb{R}\times\mathbb{R})/Heis(\mathbb{Z}\times\mathbb{Z})$ 时，它是环面上的一个 $U(1)$ 丛。

在关于弦紧化的文献中，这被称为「扭环面」（twisted torus）。

\subsection{一个自然的联络} % 原文: E.2 A Natural Connection

主 $U(1)$ 丛 (E.2) 有一个自然的联络，描述如下。给定一条路径
\begin{equation}
  p_{w_0,v} = \{w_0+tv\,|\,0\le t\le 1\}
  \tag{E.14}
  \label{eq:E.14}
\end{equation}
它在 $V$ 中从 $w_0$ 出发到 $v$，再给定 $w_0\in V$ 的一个初始提升，比如
$(w_0,z_0)\in\tilde V$，我们把 $\tilde V$ 中的平行移动点定义为
\begin{equation}
  U(p_{w_0,v}):(w_0,z_0)\to(w_0+v,\,c(v,w_0)z_0)
  \tag{E.15}
  \label{eq:E.15}
\end{equation}
沿一条闭路作平行移动——即先把 $w_0$ 平移 $v_1$，再平移 $v_2$，再平移 $-v_1$，再平移 $-v_2$——
一算（one checks）便知曲率确实就是 $\Omega$。

{\small\itshape ［原文标记：核对符号！我们要的是 $c(w_0,v)$ 吗？］}
% [原文待办标记，原文 p.187 右侧页边注：\ensuremath{\clubsuit} CHECK THE SIGN! DO WE WANT c(w_0,v)? \ensuremath{\clubsuit}]

现在，这个联络下降到环面 $V/\Lambda$ 上的主 $U(1)$ 丛 $\tilde V/\Lambda_\epsilon$ 上的一个联络。曲率仍由
$\Omega$ 给出，只是现在把 $\Omega$ 看成环面上的 $2$-形式。由此可以算出第一 Chern 类：
$(e_1,\dots,e_n)\in H^2(V/\Lambda;\mathbb{Z})$。
{\small\itshape ［原文标记：要显式证明它确实能下降（descend）到环面上。］}
% [原文待办标记，原文 p.187 右侧页边注：\ensuremath{\clubsuit} DEMONSTRATE EXPLICITLY THAT IT DOES DESCEND \ensuremath{\clubsuit}]

此外，若 $\lambda\in\Lambda$，则 $[p_{w_0,\lambda}]$ 是环面上的一条闭路，沿该路径的和乐（holonomy）为
\begin{equation}
  e^{2\pi i\Omega(\lambda,w_0)}\epsilon_\lambda^{-1}
  \tag{E.16}
  \label{eq:E.16}
\end{equation}
在把上闭链（取值在 $\mathbb{C}^*$ 中）作适当的上边缘（coboundary）改变之后，这个联络将成为下文
Hermitian 全纯线丛上的 Chern 联络。
{\small\itshape ［原文标记：这在下文要解释得更充分些。］}
% [原文待办标记，原文 p.187 右侧页边注：\ensuremath{\clubsuit} Explain this more fully below. \ensuremath{\clubsuit}]

% [Tutor Note: (E.15) 与 (E.16) 的符号要与 (E.3)、(E.4) 对齐着核对：上闭链 $c$ 的第一个自变量
% 是「被平坦移动的那一段」，而 (E.15) 里路径参数是 $v$、起点是 $w_0$，所以提升写成
% $c(v,w_0)$；但 (E.3) 的乘法约定是 $(v,z)(w_0,z_0)=(v+w_0,zz_0c(v,w_0))$，若反过来用
% $c(w_0,v)$ 则与 (E.4) 的反对称性相冲，环路的四次平移将给出 $e^{2\pi i\Omega(v_1,v_2)}$ 而非
% $e^{-2\pi i\Omega(v_1,v_2)}$——这正是页边注质疑的地方。而 (E.16) 中 $\epsilon_\lambda^{-1}$ 的
% 负幂来自 (E.6) 的等价关系（商掉右作用），若漏掉这个负幂，和乐值会相差 $\epsilon_\lambda^2$。]

\subsection{在辛环面上放一个复结构} % 原文: E.3 Putting A Complex Structure On The Symplectic Tori
{\small\itshape ［原文标记：要评论当我们有一个仿射环面时的这一点虽小、但很重要的差别。］}
% [原文待办标记，原文 p.187 右侧页边注：\ensuremath{\clubsuit} Comment on the slight, but important difference when we
% have an affine torus. \ensuremath{\clubsuit}]

我们现在想构造全纯 $\mathbb{C}^*$-丛以及与之相伴的环面上的全纯复线丛。因此需要在环面上引入一个复结构。
为此需要在 $V$ 上引入复结构 $J$：
\begin{equation}
  J:V\to V\qquad J^2=-1
  \tag{E.17}
  \label{eq:E.17}
\end{equation}
它\emph{与} $\Omega$ \emph{相容}：
\begin{equation}
  \Omega(Jv,Jw) = \Omega(v,w)
  \tag{E.18}
  \label{eq:E.18}
\end{equation}
注意这意味着
\begin{equation}
  g(v,w):=\Omega(Jv,w)=\Omega(Jw,v)
  \tag{E.19}
  \label{eq:E.19}
\end{equation}
是\emph{对称的}。

我们还将进一步假设 $g$ 正定。

现在，由于 $V$ 带有复结构，$V\otimes_{\mathbb{R}}\mathbb{C}\cong V^+\oplus V^-$ 分解为 $J$ 的 $+i$ 与
$-i$ 本征子空间。我们在 $+$ 子空间上工作，并对 $v\in V$ 定义
$v^+=\frac12(1-J\otimes i)v$。

{\small\itshape ［原文标记：我们把哪一个叫 $(1,0)$、哪一个叫 $(0,1)$，这是一种约定，在下文很要紧。］}
% [原文待办标记，原文 p.187 右侧页边注：\ensuremath{\clubsuit} Which one we call (1,0) and which (0,1) is a convention
% that matters below. \ensuremath{\clubsuit}]

\subsection{拉格朗日分解} % 原文: E.4 Lagrangian Decompositions

$V$ 的\emph{各向同性}（isotropic）子空间 $W$ 指的是这样的线性子空间：$\Omega$ 限制在其上为零，即
$\forall w,w'\in W$，$\Omega(w,w')=0$。

$(V,\Omega)$ 的拉格朗日分解指的是一个内直和
\begin{equation}
  V = W_1\oplus W_2
  \tag{E.20}
  \label{eq:E.20}
\end{equation}
其中 $W_i$ 是\emph{极大各向同性}子空间。由于 $\Omega$ 非退化，一个拉格朗日分解给出 $W_1$ 与 $W_2$
之间的非退化配对，从而给出同构 $W_2\cong W_1^{\vee}$。

上文在 (E.10) 式中引入了辛基。这样的基给出一个拉格朗日分解。对这种基引进一个记号会很方便：
对 $W_1$ 写 $\alpha_I$，$I=1,\dots,n$；对 $W_2$ 写 $\beta^J$，$J=1,\dots,n$，使得
\begin{equation}
\begin{aligned}
  \Omega(\alpha_I,\alpha_J) &= 0\\
  \Omega(\beta^I,\beta^J) &= 0\\
  \Omega(\alpha_I,\beta^J) &= e_I\delta_I{}^J
\end{aligned}
\tag{E.21}
\label{eq:E.21}
\end{equation}
若把 $v=q^I\alpha_I+p_I\beta^I$ 展开，则
\begin{equation}
  \Omega = \sum_I e_I\,dp_I\wedge dq^I
  \tag{E.22}
  \label{eq:E.22}
\end{equation}
{\small\itshape ［原文标记：需要格的整性结构，否则你大可直接取 $e_I=1$。］}
% [原文待办标记，原文 p.188 右侧页边注：\ensuremath{\clubsuit} Need integral structure of lattice otherwise you could just
% set e_I = 1. \ensuremath{\clubsuit}]
{\small\itshape ［原文标记：符号？］}
% [原文待办标记，原文 p.188 右侧页边注：\ensuremath{\clubsuit} SIGN? \ensuremath{\clubsuit}]

\subsection{环面上复结构的空间：Siegel 上半平面} % 原文: E.5 The Space Of Complex Structures On The Torus: The Siegel Upper Half-Plane

现在假设不变量 $e_I=1$。

接着选取 $(1,0)$ 型向量的一组基 $\zeta_I$。我们把 $J$ 作 $\mathbb{C}$-线性延拓到 $V_{\mathbb{C}}$，
使得按定义为
\begin{equation}
  J\cdot\zeta_I = i\zeta_I
  \tag{E.23}
  \label{eq:E.23}
\end{equation}
{\small\itshape ［原文标记：$(1,0)$ 与 $(0,1)$ 在整节里被混用了。］}
% [原文待办标记，原文 p.188 右侧页边注：\ensuremath{\clubsuit} (1,0) vs. (0,1) is all mixed up in this section. \ensuremath{\clubsuit}]

我们可以用周期矩阵（period matrix）的分量来表达复结构 $J$。后者是这样定义的：选取 $(1,0)$ 型向量的一组基
$\zeta^I$，形如
\begin{equation}
  \zeta_I := \alpha_I+\tau_{IJ}\beta^J
  \tag{E.24}
  \label{eq:E.24}
\end{equation}
{\small\itshape ［原文标记：符号？与 PiTP 2008 的约定相反。］}
% [原文待办标记，原文 p.188 右侧页边注：\ensuremath{\clubsuit} SIGN? Opposite from convention in PiTP 2008 \ensuremath{\clubsuit}]

这就给出了复向量空间之间的显式同构 $V\cong V^+$。

现在来看看二次型 $\tau_{IJ}$ 的一些性质。注意
\begin{equation}
  g(\zeta_I,\zeta_J) = \Omega(J\zeta_I,\zeta_J) = i\Omega(\zeta_I,\zeta_J)
  \tag{E.25}
  \label{eq:E.25}
\end{equation}
在 $IJ$ 上既对称又反对称，因而必须为零。于是得知 $\tau_{IJ}$ 对称，并且 $g$ 是 $(1,1)$ 型的。
注意
\begin{equation}
  g(\zeta_I,\bar\zeta_J) = 2\operatorname{Im}\tau_{IJ}
  \tag{E.26}
  \label{eq:E.26}
\end{equation}

% [Tutor Note: (E.24) 的指标位置与紧邻的正文不一致：正文说「选取 $(1,0)$ 型向量的一组基
% $\zeta^I$」（上标），公式里写的是 $\zeta_I:=\alpha_I+\tau_{IJ}\beta^J$（下标）。此处照原样转写，
% 但读者须自行统一：按 (E.24)，$\zeta$ 与 $\alpha$ 同下标、与 $\beta$ 缩并，故 $\tau_{IJ}$ 的两个
% 指标都应是「协变」的；后面 (E.30)、(E.31) 里 $Y^{-1}$、$XY^{-1}X$ 的指标位置正是按这个约定排的。
% 另外 (E.25) 中 $g(\zeta_I,\zeta_J)=i\Omega(\zeta_I,\zeta_J)$ 用的是 $g(v,w)=\Omega(Jv,w)$，
% 若有人按 $g(v,w)=\Omega(v,Jw)$ 定义，(E.25) 会整体变号，从而 (E.26) 的 $2\operatorname{Im}\tau$
% 变成 $-2\operatorname{Im}\tau$——这是本段最容易出的错。又：(E.27) 的下标排法也很特别，
% $C$ 的两个指标都写作下标（$C_{I'I}$）而 $B$ 的两个指标都写作上标（$B^{I'I}$）；要自己验证
% (E.30)，应把 (E.23) 即 $J(\alpha_I+\tau_{IJ}\beta^J)=i(\alpha_I+\tau_{IJ}\beta^J)$ 用 (E.27)
% 展开成 $\alpha,\beta$ 的线性组合再令系数相等（会用到 (E.25)--(E.26) 的 $(1,1)$ 型性质），
% 而不是直接对 $J$ 求逆。]

常与\emph{正定}度规打交道是很有用的。此时 $\operatorname{Im}\tau_{IJ}$ 是正定实二次型。

\begin{definition}
  复 $N\times N$ 对称矩阵中虚部正定的那些所构成的空间，称为\emph{Siegel 上半平面}
  $\mathcal{H}_N$。
\end{definition}

\paragraph{练习}
在实辛向量空间上给定满足 $g(v,v)>0$ 的复结构，等价于给定一个 $\tau\in\mathcal{H}_N$。

a.) 说明如何用周期矩阵来表达复结构，如下所述：复结构的作用为
\begin{equation}
\begin{aligned}
  J\cdot\alpha_I &= A^{I'}_{I}\alpha_{I'}+C_{I'I}\beta^{I'}\\
  J\cdot\beta^I &= B^{I'I}\alpha_{I'}+D^{I}_{I'}\beta^{I'}
\end{aligned}
\tag{E.27}
\label{eq:E.27}
\end{equation}
我们用下式定义向量的分量
\begin{equation}
  v = v_1^I\alpha_I+v_2^I\beta^I
  = \left(v_1^I\ v_2^I\right)\begin{pmatrix}\alpha_I\\ \beta^I\end{pmatrix}
  \tag{E.28}
  \label{eq:E.28}
\end{equation}
于是 $J$ 作用在分量上就是矩阵
\begin{equation}
  J = \begin{pmatrix} A & B\\ C & D\end{pmatrix}
  \tag{E.29}
  \label{eq:E.29}
\end{equation}
复结构的相容性蕴涵这给出一个辛矩阵。

把 (E.23) 的实部与虚部分开，并利用定义 (E.27)，我们得到 $J$ 在基 $\alpha_I,\beta^I$ 下的矩阵表达式
\begin{equation}
  J = \begin{pmatrix}
    -Y^{-1}X & Y^{-1}\\
    -Y-XY^{-1}X & XY^{-1}
  \end{pmatrix}
  \tag{E.30}
  \label{eq:E.30}
\end{equation}
其中 $\tau=X+iY$ 是 $\tau$ 的实部与虚部。可以验证 $J^2=-1$ 与 $J^{tr}\Omega J=\Omega$ 都成立。

b.) 说明在 $\alpha,\beta$ 基下度规 $g$ 为：
\begin{equation}
  g(v,w) = \left(v_1^I\ v_2^J\right)
  \begin{pmatrix}
    XY^{-1}X+Y & -XY^{-1}\\
    -Y^{-1}X & Y^{-1}
  \end{pmatrix}
  \begin{pmatrix} w_1^I\\ w_2^J\end{pmatrix}
  \tag{E.31}
  \label{eq:E.31}
\end{equation}

\paragraph{练习}
说明 $\mathcal{H}_N=Sp(2N,\mathbb{R})/U(N)$。

\paragraph{练习}
给出从整辛基到复基的变换公式是很有用的。
\begin{equation}
  \begin{pmatrix}\zeta_I\\ \bar a\zeta_I\end{pmatrix}
  = \begin{pmatrix} 1 & \tau\\ 1 & \bar a\tau\end{pmatrix}
  \begin{pmatrix}\alpha\\ \bar{\ }\end{pmatrix}
  \tag{E.32}
  \label{eq:E.32}
\end{equation}
其逆为：
\begin{equation}
  \begin{pmatrix}\alpha\\ \bar{\ }\end{pmatrix}
  = \frac{i}{2}\begin{pmatrix} \bar a\tau & -\tau\\ -1 & 1\end{pmatrix}
  Y^{-1}_{IJ}
  \begin{pmatrix}\zeta_J\\ \bar a\zeta_J\end{pmatrix}
  \tag{E.33}
  \label{eq:E.33}
\end{equation}
于是 (E.28) 的复投影为：
\begin{equation}
\begin{aligned}
  v^{(1,0)} &= -\frac{i}{2}\left(v_2^J-v_1^I\bar a\tau_{gIJ}\right)Y^{-1,JK}\zeta_K\\
  v^{(0,1)} &= \frac{i}{2}\left(v_2^J-v_1^I\tau_{IJ}\right)Y^{-1,JK}\bar a\zeta_K
\end{aligned}
\tag{E.34}
\label{eq:E.34}
\end{equation}
注意
\begin{equation}
  v = v^{(1,0)}+v^{(0,1)}
  \tag{E.35}
  \label{eq:E.35}
\end{equation}

% [Tutor Note: (E.32)--(E.34) 里反复出现的「$\bar a$」是原稿的排版残留（PDF 中印成带横线的
% 「$\bar a$」，(E.32)、(E.33) 的向量第二分量更只印出一条孤立的横线）。按上下文重构，
% 正确的写法应为：向量为 $(\alpha_I,\ \beta^I)$，矩阵为 $\binom{1\ \ \tau}{1\ \bar\tau}$，
% 即 $\zeta_I=\alpha_I+\tau_{IJ}\beta^J$、$\bar\zeta_I=\alpha_I+\bar\tau_{IJ}\beta^J$；
% 而 (E.34) 中应是 $\bar\tau_{IJ}$、$\bar\zeta_K$。用 $Y=(\tau-\bar\tau)/(2i)$ 代入
% (E.32) 的逆即可验证 (E.33)：$\binom{1\ \tau}{1\ \bar\tau}^{-1}
% =\frac{1}{\bar\tau-\tau}\binom{\bar\tau\ -\tau}{-1\ \ 1}=\frac{i}{2}\binom{\bar\tau\ -\tau}{-1\ \ 1}Y^{-1}$。
% 此处照印刷原样转写以便对照，未擅自改动。]

现在，辛群 $Sp(2N,\mathbb{R})$ 作用在复结构的空间上：$J\to hJh^{-1}$。它在周期矩阵上对应的作用
可以这样给出。

如果按向量的分量来算，那么反全纯坐标为 $v_2^I-\tau^{IJ}v_1^J$。但若分量按辛矩阵变换
\begin{equation}
  \begin{pmatrix}\tilde v^1\\ \tilde v^2\end{pmatrix}
  = \begin{pmatrix} A & B\\ C & D\end{pmatrix}
  \begin{pmatrix} v^1\\ v^2\end{pmatrix}
  \tag{E.36}
  \label{eq:E.36}
\end{equation}
那么在新的复结构下，$(Cv^1+Dv^2)-\tau(Av^1+Bv^2)$ 就是新的反全纯坐标。而它又可以写成
\begin{equation}
  (C-\tau A)v^1+(D-\tau B)v^2
  \tag{E.37}
  \label{eq:E.37}
\end{equation}
于是
\begin{equation}
  \tilde\tau = -(D-\tau B)^{-1}(C-\tau A)
  \tag{E.38}
  \label{eq:E.38}
\end{equation}

\paragraph{练习（$Sp(2N,\mathbb{R})$ 的自同构）}
变换律有许多种形式，因为我们既可以用被动（passive）变换也可以用主动（active）变换，还可以与各种
自同构复合。

辛群由 $g^{tr}\mathcal{J}g=\mathcal{J}$ 定义，因此
{\small\itshape ［原文标记：核对！！］}
% [原文待办标记，原文 p.191 右侧页边注：\ensuremath{\clubsuit} CHECK!! \ensuremath{\clubsuit}]
\begin{equation}
  \begin{pmatrix} A & B\\ C & D\end{pmatrix}
  \to
  \begin{pmatrix} D & -C\\ -B & A\end{pmatrix}
  \tag{E.39}
  \label{eq:E.39}
\end{equation}
是一个自同构，于是复合之后得到
\begin{equation}
  \tilde\tau = (A+\tau C)^{-1}(B+\tau D)
  \tag{E.40}
  \label{eq:E.40}
\end{equation}
再借助别的自同构，说明我们可以把它写成
\begin{equation}
  \tilde\tau = (A\tau+B)(C\tau+D)^{-1}
  \tag{E.41}
  \label{eq:E.41}
\end{equation}
并证明 $\tilde\tau$ 是对称的。

现在，如果我们关心的是环面上与某个辛结构相容的复结构，那么我们希望辛变换保持格不变。覆盖空间上
复结构的空间由 $\mathcal{H}_N\cong Sp(2N,\mathbb{R})/U(N)$ 给出。但我们还应把仅仅相差格的一组基
变换的复结构等同起来：

\begin{theorem}
  $\mathbb{R}^{2N}/\mathbb{Z}^{2N}$ 上复结构的空间是
  \begin{equation}
    Sp(2N,\mathbb{Z})\backslash\mathcal{H}_N
    \cong Sp(2N,\mathbb{Z})\backslash Sp(2N,\mathbb{R})/U(N)
    \tag{E.42}
    \label{eq:E.42}
  \end{equation}
\end{theorem}

注：可与 $\mathbb{R}^{2N}$ 上与某个欧氏结构相容的复结构空间作比较，后者是 $O(2N)/U(N)$。

最后，若作整值辛变换，我们只是改变了拉格朗日分解的选取。但线丛及其截面空间是环面内蕴的。这蕴涵
theta 函数具有值得注意的变换性质。

（整值）辛群由三类变换生成：
\begin{equation}
  \begin{pmatrix} A & 0\\ 0 & A^{tr,-1}\end{pmatrix}
  \tag{E.43}
  \label{eq:E.43}
\end{equation}
它把 $p$ 与 $q$ 各自重新排列。
\begin{equation}
  \begin{pmatrix} 1 & B\\ 0 & 1\end{pmatrix}
  \tag{E.44}
  \label{eq:E.44}
\end{equation}
它把一些 $p$ 加到 $q$ 上。
\begin{equation}
  \begin{pmatrix} 0 & -1_N\\ 1_N & 0\end{pmatrix}
  \tag{E.45}
  \label{eq:E.45}
\end{equation}
它交换 $p$ 与 $q$。

关于这三类变换确实生成整值辛群的证明，见 D. Mumford，\textit{Tata Lectures on Theta}，第 I 卷，
p.202。
